\documentclass[conference]{IEEEtran}
\IEEEoverridecommandlockouts

\usepackage{cite}
\usepackage{amsmath,amssymb,amsfonts}
\usepackage{algorithmic}
\usepackage{graphicx}
\usepackage{textcomp}
\usepackage{xcolor}
\usepackage{booktabs}

\begin{document}

\title{Blockchain-Based IoT Device Authentication Using PUF and HMAC-SHA256}

\author{\IEEEauthorblockN{Dundi Mahendar Reddy, GangaRamPrasad Kotakonda, Rahul Sambaturu, Asif Shaik}
\IEEEauthorblockA{Department of Computer Science and Engineering\\
SRM University AP\\
Andhra Pradesh, India\\
Project ID: CAPSTONE 2022 24183 4}
}

\maketitle

\begin{abstract}
This paper presents a hybrid authentication framework for Internet of Things (IoT) devices that combines Physical Unclonable Function (PUF)-derived secrets, HMAC-SHA256 message authentication, and Hyperledger Fabric blockchain logging. The framework addresses device spoofing, tampering, and replay attacks while preserving low-latency gateway response behavior through asynchronous blockchain logging. We implement a complete device-to-gateway-to-blockchain pipeline and validate it using integrated security test cases including valid authentication, replay attempts, invalid signature submissions, and unregistered device access. Our prototype demonstrates practical deployment with successful chaincode lifecycle completion and immutable ledger evidence generation. A benchmark run demonstrates 100\% successful authentications under controlled load with a measured median latency of 27.59 ms and throughput of 456.20 authentications per second for 200 requests at concurrency 20. The framework is reproducible via one-command automation scripts and is suitable as a foundation for production hardening and future biometric expansion.
\end{abstract}

\begin{IEEEkeywords}
IoT security, PUF, HMAC-SHA256, Hyperledger Fabric, replay attack prevention, blockchain auditability
\end{IEEEkeywords}

\section{Introduction}
IoT deployments require lightweight but secure authentication primitives due to resource constraints and high attack exposure. Static credentials and centralized logs create risks including credential cloning, tampering, and non-repudiation gaps. This work proposes a practical framework integrating PUF-based device identity, HMAC verification, and blockchain-backed audit trails.

\subsection{Contributions}
The main contributions are:
\begin{enumerate}
\item End-to-end IoT authentication pipeline across device, gateway, and blockchain layers.
\item Replay resistance via timestamp freshness and nonce reuse checks.
\item Asynchronous Fabric logging to avoid authentication path blocking.
\item Reproducible deployment, testing, and benchmark automation.
\end{enumerate}

\section{System Design}
\subsection{Device Layer}
A Python client derives a deterministic PUF-like secret and computes HMAC-SHA256 over payload fields $(message, timestamp, nonce)$.

\subsection{Gateway Layer}
An Express gateway validates request schema, device registration, timestamp freshness window ($\pm 60$ s), nonce uniqueness, and HMAC correctness. Constant-time comparison is used to reduce timing leakage.

\subsection{Blockchain Layer}
Successful authentications are logged to Hyperledger Fabric on channel \texttt{mychannel}. Chaincode deployment follows install, approve, commit, and query lifecycle.

\section{Implementation}
The implementation consists of:
\begin{itemize}
\item \textbf{Device:} Python modules for PUF simulation and HMAC token generation.
\item \textbf{Gateway:} Node.js/Express API endpoint \texttt{/auth} with security checks.
\item \textbf{Fabric Ops:} Automated scripts for network setup, chaincode deployment, and ledger inspection.
\item \textbf{Benchmarking:} Scripted load generation with latency percentile and throughput output.
\end{itemize}

\section{Experimental Results}
\subsection{Security Behavior}
The demo test suite validated expected decisions:
\begin{itemize}
\item Valid request: accepted
\item Replay request: rejected
\item Invalid HMAC: rejected
\item Unregistered device: rejected
\item Invalid payload: rejected
\end{itemize}

\subsection{Performance Snapshot}
A benchmark run produced the metrics in Table~\ref{tab:perf}.

\begin{table}[htbp]
\caption{Authentication Performance Snapshot}
\label{tab:perf}
\centering
\begin{tabular}{lr}
\toprule
Metric & Value \\
\midrule
Requests & 200 \\
Concurrency & 20 \\
Success Rate & 100\% \\
Throughput & 456.20 auth/s \\
p50 Latency & 27.59 ms \\
p95 Latency & 92.49 ms \\
p99 Latency & 102.79 ms \\
\bottomrule
\end{tabular}
\end{table}

\section{Discussion}
The architecture achieves strong functional security coverage and practical observability via immutable logging. While median latency is low, tail latency (p95/p99) requires optimization under higher concurrency. Future work should include repeated trials, confidence intervals, and multi-node gateway scaling.

\section{Conclusion}
This work demonstrates a deployable and reproducible blockchain-backed IoT authentication framework combining PUF and HMAC-SHA256. The integrated pipeline, security behavior, and automation artifacts support both capstone delivery and transition to publication-grade evaluation.

\section*{Acknowledgment}
The authors thank Dr. Mallavalli Sitharam for supervision and technical guidance.

\begin{thebibliography}{00}
\bibitem{b1} Hyperledger Fabric Documentation, "Hyperledger Fabric," https://hyperledger-fabric.readthedocs.io/.
\bibitem{b2} H. Krawczyk, M. Bellare, and R. Canetti, "HMAC: Keyed-Hashing for Message Authentication," RFC 2104.
\bibitem{b3} C. Herder, M.-D. Yu, F. Koushanfar, and S. Devadas, "Physical Unclonable Functions and Applications: A Tutorial," Proc. IEEE, 2014.
\bibitem{b4} NIST, "Considerations for Managing IoT Cybersecurity and Privacy Risks," NISTIR 8228.
\end{thebibliography}

\end{document}
\...
